Para resolver esta falta de acompañamiento sin caer en el extremo de resolverle el trabajo al alumno, la pedagogía moderna propone estrategias de instrucción guiada, destacando entre ellas el método socrático \cite{kargupta2024instruct}. Este enfoque no busca dar respuestas, sino establecer un diálogo donde el tutor guía al estudiante mediante preguntas pequeñas y estratégicas. El objetivo es que el alumno "tire del hilo" de su propio razonamiento hasta descubrir, por sí mismo, dónde está el fallo en su lógica \cite{alhossami2024can}. Al obligar al estudiante a verbalizar y reflexionar sobre lo que está intentando hacer, se consigue un aprendizaje mucho más profundo que el que se obtiene con una corrección automática \cite{aleven2002effective}.

Este apoyo, conocido como \textit{scaffolding}, permite que el alumno no se rinda ante el primer bloqueo, ya que la IA actúa como un soporte temporal que se adapta a su nivel \cite{wood1976role}. En el contexto actual de la Inteligencia Artificial, este método se vuelve una pieza fundamental para garantizar que la tecnología sea una ayuda y no una herramienta de plagio. Mediante una cuidada ingeniería de \textit{prompts}, podemos configurar sistemas que emulen la paciencia de un tutor humano, combatiendo así el riesgo de que el alumno delegue su esfuerzo cognitivo en la máquina y pierda la oportunidad de aprender realmente a programar \cite{chen2024plagiarism}.